% ---------------------------------------------------------------------------
% V20 manuscript — SHARED BODY (science identical in blind and supervisor
% variants). Compiled from paper/v20-final on branch paper/v20-final.
% Target: 6 pages preferred, 8 pages hard maximum (IEEE conference format).
% ---------------------------------------------------------------------------
\maketitle

\begin{abstract}
Repository-level language models must first decide which files a requested
change affects, but the explicit impact plan that encodes this decision is
itself an output-budget consumer: on a 144-file candidate universe an explicit
per-file policy serializes roughly 140--150 decision records per task. We study
\textit{Preserve-by-Omission}, a sparse representation in which only
non-default decisions are emitted and every omitted candidate deterministically
decodes to ``preserve.'' A controlled 16K-cap encoding ablation on six curated
django CMS mechanism tasks shows the sparse form uses about 90\% fewer
completion tokens and 96\% fewer serialized records at equal operational
validity, while selection fidelity is scenario-dependent rather than
universally better. On ten independent, previously unseen real historical
django CMS changes (sixty cells, two arms, three nested repetitions), the
output-cost advantage persists (mean completion 8,446 vs.\ 599 tokens; total
cost \$0.298 vs.\ \$0.062), while paired task-level F1 differences are small
and cross zero. A compact system-level comparison against the LocAgent
graph-guided localization framework (pinned upstream commit
\texttt{4935b557326c154bad8e8dcf3747cc8d32d1f387}, OpenRouter-routed
Qwen3-Coder) on the same ten tasks shows comparable-to-lower file-level F1 at
roughly two orders of magnitude more inference cost, with a 50\%
empty/non-usable localization rate (two timeouts, one context-length failure,
two completed-but-empty outcomes). Sparse impact plans offer a large and
reproducible output/cost reduction for repository-level impact selection;
semantic fidelity and omission risk remain central open problems.
\end{abstract}

\section{Introduction}
\IEEEPARstart{R}{epository}-level large language models (LLMs) increasingly
drive software evolution tasks such as issue resolution and code repair
\cite{jimenez2024swebench,yang2024sweagent}. Before any edit can be written, a
model must decide \emph{which repository-relative files are affected by the
requested change}. We call this decision the \emph{impact-selection stage}. It
is a prerequisite for faithful editing: an agent that never locates a changed
file cannot fix it, and an agent that proposes an excessively large scope wastes
output budget and risks touching unrelated code.

Modern approaches to this stage fall into two broad families. Agentic systems
\cite{chen2025locagent,yang2024sweagent,xia2024agentless} navigate the
repository interactively, issuing many model calls per task; retrieval-based
systems \cite{zhang2023repocoder,tian2024repobench} rank candidate files by
similarity or graph proximity. Both families measure \emph{selection
effectiveness} (how often the correct files are chosen) but rarely account for
the \emph{serialization cost} of expressing the chosen scope. For an explicit
impact-plan representation over a large candidate universe, this cost can
dominate the interaction budget.

This paper isolates a single, separable problem: \emph{given the same decision
semantics, how much output budget does the representation of an impact plan
consume, and does the cheaper representation preserve selection fidelity?} We
contribute (1) Preserve-by-Omission, a sparse explicit plan in which omitted
candidates deterministically decode to ``preserve''; (2) a controlled
16K-cap encoding ablation on six curated django CMS mechanism tasks; (3) a
real-commit replication on ten independent historical django CMS changes; and
(4) a compact system-level comparison against the LocAgent graph-guided
localization framework on the same ten tasks. We find the output-cost
advantage is large and reproducible, while semantic fidelity is
heterogeneous and omission risk (high false-negative rate) remains a central
weakness.

\section{Related Work}
\IEEEPARstart{S}{oftware}-repository benchmarks such as SWE-bench
\cite{jimenez2024swebench} evaluate full issue resolution end-to-end. Within
that pipeline, localization is often implicit: an agent searches the
repository until it finds an edit location. SWE-agent \cite{yang2024sweagent}
explicitly instruments a file editor, while Agentless \cite{xia2024agentless}
shows that a lightweight locate-then-repair pipeline can be competitive without
an interactive loop.

For the localization step itself, RepoCoder \cite{zhang2023repocoder} and
RepoBench \cite{tian2024repobench} frame file/line completion as ranked
retrieval. CodePlan \cite{bairi2023codeplan} plans repository changes as
dependency-constrained edit chains. Most relevant to our comparison, LocAgent
\cite{chen2025locagent} is a graph-guided agentic framework for code
localization: it parses a repository into a directed heterogeneous graph, lets
an LLM agent traverse the graph and BM25 index, and reports ranked
files/entities. Its published file-level accuracy (\textasciitilde92.7\%
Acc@5) is reported on its own fine-tuned checkpoint and its own localization
benchmark.

Our work differs in two ways. First, we measure the \emph{representation
cost} of an explicit impact policy, which the localization literature does not
isolate: we hold decision semantics constant and vary only serialization. A
full policy over a 144-path universe requires roughly 140--152 explicit
records; Preserve-by-Omission reduces this to a handful. Second, we compare
systems on the \emph{same} ten real held-out tasks under a shared common
evaluator, and we explicitly do not reproduce LocAgent's published
fine-tuned configuration: our P5 baseline runs the pinned upstream framework
with an OpenRouter-routed Qwen3-Coder model, which makes the comparison
\emph{system-level}, not an algorithm ablation.

\section{Method: Preserve-by-Omission}
\IEEEPARstart{T}{he} input to impact selection is a repository, a candidate
universe $U$ of repository-relative paths, and a natural-language change
request. An \emph{impact plan} assigns each candidate exactly one action from a
fixed vocabulary; in this study the vocabulary is
\{PRESERVE, REGENERATE, VALIDATE, HUMAN\_REVIEW\}. A \emph{full} explicit plan
serializes one decision per candidate. Because most candidates are unaffected
by any single change, the full serialization is dominated by hundreds of
repeated ``PRESERVE'' rows.

\textbf{Preserve-by-Omission.} The sparse representation emits only
non-PRESERVE decisions; every omitted candidate is reconstructed
deterministically as PRESERVE. The decoding function $D_s$ maps the emitted
decision set $E_s(\pi)$ back to a complete policy: $D_s(E_s(\pi)) = \pi$ for
every valid policy. This gives an exact, lossless encoding of the decision
semantics and separates \emph{representation cost} from \emph{decision
quality}: both arms carry identical action semantics and differ only in how the
plan is serialized.

\textbf{Controlled treatment.} We reuse one common JSON-schema prompt template
and a frozen candidate-ID map. The only difference between the arms is the
serialization-policy block: \emph{Full-v2} instructs explicit decisions for all
candidates; \emph{Sparse-v2} instructs explicit non-PRESERVE decisions only. A
frozen representation-equivalence check proves $D_s(E_s(\pi)) = \pi$ on all
valid outputs, and a prompt-control check proves the two rendered prompts are
byte-identical after stripping the policy block. The model is
Qwen3-Coder-480B-A35B-Instruct \cite{qwen2025qwen3} via OpenRouter with
temperature 0, structured-output JSON schema, and a 16,384-token completion cap
for both arms.

\begin{figure}[t]
\centering
\fbox{\parbox{0.94\linewidth}{\centering
\begin{minipage}{0.94\linewidth}
\raggedright
\small
\textbf{Full-v2:} ``For each candidate $1\ldots144$, emit one decision.''\\
\emph{Output:} 144 rows (mostly PRESERVE), $\approx$8.4k tokens.\\[2pt]
\textbf{Sparse-v2:} ``Emit only non-PRESERVE decisions.''\\
\emph{Output:} $\approx$4 rows, $\approx$0.6k tokens.\\[2pt]
\textbf{Decode:} $D_s$ fills every omitted candidate with PRESERVE\\
$\Rightarrow$ identical complete policy, exact reconstruction.
}}
}
\caption{Preserve-by-Omission: both arms carry identical decision semantics;
they differ only in serialization. The sparse form emits a handful of rows and
the decoder reconstructs the complete policy deterministically.}
\label{fig:pbo}
\end{figure}

\section{Experimental Design}
\IEEEPARstart{W}{e} run three complementary studies, all on the django CMS 5.0.0
repository and all scored against the same kind of reference. The historical
set of changed production files for a real commit is an \emph{observed
change-set proxy}: it is evidence about what was actually edited, not a perfect
semantic ground truth. Selection metrics (precision, recall, F1, FNR) compare
each predicted file set against this proxy. The reference is applied
evaluation-only and never appears in any prompt.

\textbf{Real-commit corpus (M4A).} A deterministic miner over django CMS
history produced 40 clean scientific changes from the newest 6,000 ancestors
(\textasciitilde2016--2025), with a parent-only public/hidden split and a
leakage barrier that rejects any case whose visible intent already names the
changed paths. Ninety-two leaked candidates were excluded, exact-duplicate and
related-change removal ran on the remainder, and a metadata-only split freeze
(TRAIN 24 / VALIDATION 6 / HELD\_OUT\_TEST 10, seed frozen before any model
output) fixed the ten held-out tasks used by P1 and P5.

\textbf{M1 — controlled mechanism.} Six curated development/mechanism
scenarios (002, 004, 005, 006, 007, 008), two arms $\times$ five repetitions
= 60 cells, completion cap 16,384. This is a controlled encoding ablation: the
completion cap is non-binding for both arms (0 truncations), so any output
difference isolates serialization, not truncation. The six scenarios are
independent task units; the five repetitions are nested observations (no
$n=30$ independence claim). A separate 4,096-cap feasibility boundary (M1A)
is reported only as mechanism evidence: at that cap the Full arm truncates
before emitting all required decisions while Sparse completes, which is why
all semantic comparisons in this paper use the non-binding 16K cap (avoiding
survivor bias toward the sparse arm).

\textbf{P1 — real-commit replication.} Ten independent held-out real django CMS
commits (HELD\_OUT\_TEST), two arms $\times$ three nested repetitions = 60
cells, same 16K cap, same prompts except the serialization block. Paired
task-level deltas use a 10,000-resample bootstrap over the ten independent
tasks, seeded and frozen; nested repetitions are pooled per task, never
treated as independent examples.

\textbf{P5 — system-level LocAgent baseline.} The same ten held-out tasks are
run once each through the LocAgent framework \cite{chen2025locagent} at pinned
upstream commit \texttt{4935b557326c154bad8e8dcf3747cc8d32d1f387}, using
OpenRouter-routed Qwen3-Coder (the framework's upstream protocol uses
temperature 1, MRR ranking, a 900~s per-attempt timeout, and one execution per
task). This is a \emph{framework baseline}, not a reproduction of the
published fine-tuned LocAgent configuration, and it is evaluated with the
same common evaluator against the same observed change-set proxy. P5 cost is a
normalized estimate under a frozen pricing snapshot because the backend
provider was OpenRouter-routed and not proven per call.

\subsection{Validity design}
Six curated scenarios are development/mechanism evidence, not held-out
confirmation; only the ten P1/P5 tasks are independent held-out data. The
common evaluator applies one fixed policy (exact emitted file set), and the
LocAgent-native ranking metric (Acc@K) is kept separate from the common F1
columns because the two measure different things. A per-case usage ledger
records model calls, tokens, latency, and estimated cost for P5; the model-call
count is the authoritative ledger count, never the number of persisted output
files.

\subsection{Reproducibility and auditing}
Every study passes six frozen pre-benchmark validation gates (dataset, prompt,
pipeline smoke, dry-run, integration, metric) before any scientific call, and
an independent zero-API audit after. All raw evidence is hash-pinned: P1's 60
per-cell raw responses carry SHA-256 sidecars; the P5 per-call usage ledger and
merged ranked outputs are immutable. A strengthened zero-API audit (28 checks)
verifies the common P/R/F1/FNR arithmetic from raw files, the official native
Acc@K semantics, the empty-outcome failure taxonomy against raw logs, the
provider-route wording, and the consistent efficiency denominator. These
checks exist so that the reporting errors this paper corrects (item-hit counts
labelled as task accuracy, a ``5/10 timeout'' taxonomy, and mixed validity
denominators) fail closed in future releases.

\section{Results}
\subsection{M1 — Controlled mechanism}
At the 16K cap both arms achieve 60/60 operational validity with zero
truncations, so the representation is the only systematic difference.

\begin{table}[t]
\caption{M1 controlled 16K encoding ablation (6 curated scenarios $\times$ 2
arms $\times$ 5 reps = 60 cells; all valid).}
\label{tab:m1}
\centering
\small
\begin{tabular}{lcc}
\toprule
\textbf{Metric} & \textbf{Full-v2} & \textbf{Sparse-v2} \\
\midrule
Valid / cells & 30/30 & 30/30 \\
Truncations & 0 & 0 \\
Mean completion tokens & 8,383 & 809 \\
Mean serialized records & 144.0 & 5.9 \\
Precision & 0.452 & 0.721 \\
Recall & 0.775 & 0.883 \\
F1 & 0.571 & 0.794 \\
Total API cost (USD) & 0.275 & 0.048 \\
\bottomrule
\end{tabular}
\end{table}

Table~\ref{tab:m1} shows the controlled effect: Sparse-v2 reduces mean
completion output by about 90\% (8,383$\rightarrow$809) and serialized records
by about 96\% (144.0$\rightarrow$5.9) at identical validity and no truncation.
Selection fidelity improves on four of six scenarios and degrades on two;
scenario 006 is the counterexample (Sparse recall 0.333 vs.\ Full 0.800). We
therefore describe the mechanism result as a controlled
\textit{serialization-cost} effect, not a universal semantic improvement.

\subsection{P1 — Real-commit replication}
The ten held-out real commits are substantially harder and noisier than the
six curated mechanism tasks. Both arms still complete 30/30 valid cells per arm
with zero truncations.

\begin{table}[t]
\caption{P1 real-commit held-out evaluation (10 tasks $\times$ 2 arms $\times$
3 nested reps = 60 cells; all valid).}
\label{tab:p1}
\centering
\small
\begin{tabular}{lcc}
\toprule
\textbf{Metric} & \textbf{Full-v2} & \textbf{Sparse-v2} \\
\midrule
Valid / cells & 30/30 & 30/30 \\
Precision & 0.339 & 0.387 \\
Recall & 0.369 & 0.261 \\
F1 & 0.353 & 0.312 \\
FNR & 0.631 & 0.739 \\
Mean completion tokens & 8,445.8 & 599.0 \\
Mean serialized records & 144.0 & 4.07 \\
Total API cost (USD) & 0.2976 & 0.0623 \\
\bottomrule
\end{tabular}
\end{table}

As Table~\ref{tab:p1} shows, the output-cost advantage transfers to unseen
historical changes: Sparse-v2 cuts mean completion tokens by 7,848
(bootstrap 95\% CI [--8,136, --7,597], excludes zero), serialized records by
139.9 [--143.3, --137.0], and per-task cost by \$0.0235 (CI excludes zero).
Task-level $\Delta$F1 (Sparse $-$ Full) is $-$0.009 with 95\% CI
[$-$0.130, $+$0.119], which crosses zero. Both arms have high false-negative
rates (0.63 and 0.74); on these real changes omission risk is a central
weakness for both representations.

\begin{table}[t]
\caption{P1 per-task F1 (pooled over three nested repetitions per task).
Fidelity is heterogeneous: the arms agree on easy and hard tasks but trade off
precision and recall on the middle difficulty.}
\label{tab:p1tasks}
\centering
\scriptsize
\setlength{\tabcolsep}{3pt}
\begin{tabular}{lcc}
\toprule
\textbf{Task} & \textbf{Full F1} & \textbf{Sparse F1} \\
\midrule
4307e1b8c2e2 & 0.500 & 0.400 \\
50c3576080be & 0.333 & 0.143 \\
630a50361ada & 0.167 & 0.000 \\
66c70394c9e1 & 0.214 & 0.545 \\
75978fb1c3ad & 1.000 & 1.000 \\
8d50660e7bcf & 0.929 & 0.571 \\
9e33db4f4660 & 0.118 & 0.400 \\
b39799f9fc1c & 0.341 & 0.357 \\
ba16eb9a1d09 & 0.000 & 0.000 \\
fdda30c271f0 & 0.267 & 0.375 \\
\bottomrule
\end{tabular}
\end{table}

Table~\ref{tab:p1tasks} shows the per-task F1 values that drive the
inconclusive paired delta. On the two easy tasks both arms reach F1 1.0; on
the two hardest tasks both arms score 0; on the middle six tasks the arms
trade off precision and recall, and neither arm dominates. This is consistent
with a representation that changes the output shape (and its cost) without
changing the intrinsic difficulty of each task, and it is why we report a
controlled cost effect rather than a semantic superiority claim.

\subsection{P5 — LocAgent framework baseline}
Table~\ref{tab:p5} reports the system-level comparison on the same ten tasks.
Execution/validity denominators are explicit: P1 arms ran 10 tasks $\times$ 3
nested repetitions = 30 cells each; LocAgent ran 10 tasks $\times$ 1 execution
= 10 outcomes.

\begin{table}[t]
\caption{P5 system-level shared-task comparison on the same ten held-out
commits. LocAgent: pinned upstream
\texttt{4935b557326c154bad8e8dcf3747cc8d32d1f387}, OpenRouter-routed
Qwen3-Coder, temperature 1, MRR, 900~s timeout, one execution per task.}
\label{tab:p5}
\centering
\scriptsize
\setlength{\tabcolsep}{3pt}
\begin{tabular}{lcccccc}
\toprule
\textbf{System} & \textbf{Tasks} & \textbf{Runs} & \textbf{Non-empty} &
\textbf{Empty} & \textbf{P} & \textbf{F1} \\
\midrule
Full-v2 & 10 & 30 & 30 & 0 & 0.339 & 0.353 \\
Sparse-v2 & 10 & 30 & 30 & 0 & 0.387 & 0.312 \\
LocAgent & 10 & 10 & 5 & 5 & 0.435 & 0.333 \\
\bottomrule
\end{tabular}
\end{table}

LocAgent produced a usable file set on 5 of 10 tasks. The other five are
fail-closed empty and are \emph{not} all timeouts: the raw logs show two 900~s
timeouts, one context-length \texttt{BadRequestError}, and two
completed-but-empty outcomes (the upstream flow logged ``succeed'' with no
parseable file set) --- a 50\% empty/non-usable localization rate, not a 50\%
timeout rate.

\begin{table}[t]
\caption{Taxonomy of the five empty LocAgent outcomes, classified from the raw
execution logs. All five are persisted fail-closed, but only two are timeouts.}
\label{tab:p5fail}
\centering
\scriptsize
\setlength{\tabcolsep}{3pt}
\begin{tabular}{lll}
\toprule
\textbf{Task} & \textbf{Class} & \textbf{Log evidence} \\
\midrule
4307e1b8c2e2 & 900~s timeout & ``exceeded timeout'' \\
66c70394c9e1 & context-length & ``maximum context\\
 & BadRequest & length'' \\
9e33db4f4660 & completed-but-empty & ``succeed'', no files \\
b39799f9fc1c & completed-but-empty & ``succeed'', no files \\
fdda30c271f0 & 900~s timeout & ``exceeded timeout'' \\
\bottomrule
\end{tabular}
\end{table}

Pooled micro F1 for LocAgent (0.333) sits between Sparse (0.312)
and Full (0.353); paired task-level $\Delta$F1 bootstrap CIs cross zero for
both comparisons (LocAgent $-$ Full: $-$0.068 [$-$0.250, $+$0.170]; LocAgent
$-$ Sparse: $-$0.061 [$-$0.306, $+$0.241]). LocAgent's native ranking on the
tasks that complete is Acc@1 4/10, Acc@3 4/10, Acc@5 2/10 under the official
LocAgent metric; these are reported separately and are not comparable to our
F1. The failure taxonomy in Table~\ref{tab:p5fail} is reported so that the
robustness profile is not conflated with a single timeout mechanism.

The efficiency contrast is large and is stated with one consistent denominator
(mean per execution/task):

\begin{table}[t]
\caption{Mean per-execution/task inference efficiency on the same ten tasks.
One consistent denominator: P1 means are per cell ($n=30$), LocAgent means are
per task execution ($n=10$). P5 cost is a normalized estimate under a frozen
pricing snapshot.}
\label{tab:efficiency}
\centering
\scriptsize
\setlength{\tabcolsep}{3pt}
\begin{tabular}{lrrr}
\toprule
\textbf{System} & \textbf{Mean total} & \textbf{Mean cost} & \textbf{Calls} \\
 & \textbf{tokens} & \textbf{(USD)} & \\
\midrule
Full-v2 & 13,362 & 0.00992 & 30 \\
Sparse-v2 & 5,530 & 0.00208 & 30 \\
LocAgent & 3,283,177 & 0.99288 & 402 \\
\bottomrule
\end{tabular}
\end{table}

As Table~\ref{tab:efficiency} shows, LocAgent uses about 246$\times$ the mean
tokens and 100$\times$ the mean cost of Full-v2 (and about 594$\times$ /
478$\times$ of Sparse-v2), totaling 402 model calls and roughly \$9.93 of
normalized estimated cost (32.8M tokens) versus \$0.30 (Full) and \$0.06
(Sparse). On these ten tasks the accuracy--inference-cost frontier does not
favor the graph-agent baseline; the cost advantage of Preserve-by-Omission is
not offset by superior localization accuracy.

\section{Discussion}
\IEEEPARstart{T}{he} evidence supports three claims. First, the representation
of an impact plan is a first-class output-budget cost: on a 144-path universe
an explicit full policy costs roughly an order of magnitude more output than a
sparse one, and this holds in a controlled ablation with a non-binding cap.
Second, this cost advantage transfers to independent real historical changes:
Sparse-v2 reduces completion tokens, serialized records, and cost with
bootstrap CIs excluding zero, at small and statistically inconclusive F1
deltas. Third, at the system level, an interactive graph-guided agent does not
deliver a clear accuracy advantage on these ten tasks at its much higher
inference cost.

These findings have a design implication for repository-level editing. The
high false-negative rate of both our arms (0.63--0.74) means the cheap sparse
plan should not be trusted as final scope. A cost-effective architecture would
use a sparse first-pass plan, flag omission-risk candidates, and escalate only
those for bounded verification --- rather than running an always-on
high-cost agent on every task. LocAgent's graph traversal suggests dependency
proximity as a useful soft signal for that escalation, but no graph-assisted
inference was executed in this study, so we make no causal claim about graph
benefit. The P5 result also illustrates a system-level trade-off that is
independent of the sparse representation: an interactive localization agent can
spend two orders of magnitude more inference budget without a corresponding
accuracy gain, which strengthens the case for budget-aware selection
strategies.

\section{Threats to Validity}
\IEEEPARstart{W}{e} consolidate the main threats. (1) The observed
change-set proxy is not perfect semantic ground truth; the historical diff may
under- or over-state the files a model ``should'' flag. (2) The ten held-out
tasks come from a single repository (django CMS); external generalization to
other repositories is not measured. (3) M1 uses six curated scenarios as
independent units; five repetitions are nested, so $n=6$ is the inferential
unit for the controlled study. (4) Semantic effects are heterogeneous, with
scenario 006 degrading under the sparse arm; no universal superiority claim is
made. (5) The P1 $\Delta$F1 CIs cross zero, so we do not claim a detected
difference --- and crossing zero is not evidence of equivalence either. (6) The
P5 comparison is system-level, not an algorithm ablation: protocols,
temperature, interaction budget, and model serving differ; LocAgent's backend
provider was OpenRouter-routed and not proven per call, and its cost is a
normalized estimate under a frozen pricing snapshot. (7) The 4,096-cap
feasibility boundary is not used for semantic comparison because it censors
the Full arm (survivor bias); all semantic comparisons use the non-binding 16K
cap. (8) High FNR is a weakness, not a success: both representations omit a
substantial fraction of observed proxy files.

\section{Conclusion}
\IEEEPARstart{W}{e} introduced Preserve-by-Omission, a sparse explicit impact
plan that deterministically reconstructs a complete policy from only its
non-default decisions. In a controlled 16K-cap encoding ablation it reduces
completion output by about 90\% and serialized records by about 96\% at equal
operational validity; on ten independent real historical django CMS commits the
output/cost advantage persists while semantic fidelity remains heterogeneous
and paired F1 differences are statistically inconclusive. A compact
system-level comparison shows the LocAgent graph-guided framework does not
outperform the cheap sparse planner on accuracy while costing roughly two
orders of magnitude more inference. The central open problem is omission risk:
high false-negative rates motivate selective, bounded verification of
low-cost impact plans rather than always-on agentic search. All experiments,
scores, and audits are reproducible from frozen raw evidence in the public
repository.

\begin{thebibliography}{9}
\bibitem{chen2025locagent} Z.~Chen et al., ``LocAgent: Graph-guided LLM agents
for code localization,'' in \emph{Proc.\ ACL (Long Papers)}, 2025.
\bibitem{jimenez2024swebench} C.~E. Jimenez et al., ``SWE-bench: Can language
models resolve real-world GitHub issues?'' in \emph{Proc.\ ICLR}, 2024.
\bibitem{xia2024agentless} C.~S. Xia et al., ``Agentless: Demystifying
LLM-based software engineering agents,'' \emph{arXiv:2407.01489}, 2024.
\bibitem{yang2024sweagent} J.~Yang et al., ``SWE-agent: Agent-computer
interfaces enable automated software engineering,'' in \emph{Proc.\ NeurIPS},
2024.
\bibitem{zhang2023repocoder} F.~Zhang et al., ``RepoCoder: Repository-level
code completion through iterative retrieval and generation,'' in
\emph{Proc.\ EMNLP}, 2023.
\bibitem{tian2024repobench} R.~Tian et al., ``RepoBench: Benchmarking
repository-level code auto-completion systems,'' in \emph{Proc.\ ICLR}, 2024.
\bibitem{bairi2023codeplan} R.~Bairi et al., ``CodePlan: Repository-level
coding using LLMs and planning,'' in \emph{Proc.\ NeurIPS}, 2023.
\bibitem{qwen2025qwen3} Qwen Team, ``Qwen3 technical report,''
\emph{arXiv:2505.09388}, 2025.
\end{thebibliography}